\section{引言}
\subsection{研究背景}
临床自然语言处理的核心目标之一，是从电子病历、住院记录、出院小结和专科会诊文书中提取能够支持医疗分析与辅助决策的高价值信息。与结构化检验指标不同，临床文本保留了医生对病情的叙述、判断和推理过程，因此在表型识别、疾病风险预测、药物关系抽取、病历质控和去标识化处理等任务中具有不可替代的作用。然而，临床文本同时又具有篇幅长、缩写多、术语密集、叙事结构不稳定以及标签边界模糊等特点，使得其建模难度明显高于通用文本。

近年来，预训练语言模型特别是领域特化编码器显著提升了医学文本表示能力。本文采用的 \encoder{} 即是一种面向生物医学和临床场景的长上下文编码器，其目标是在更长文本范围内保留医学语义依赖关系 \citep{sounack2025bioclinicalmodernbertstateoftheartlongcontext}。这类模型为临床 NLP 提供了强有力的上游表示基础，也使研究重点逐渐从“如何获得更好的文本表示”扩展到“如何基于强表示完成更可靠的下游决策”。

但在真实医疗应用中，仅有更强的编码表示仍然不够。临床任务不仅关心模型能否输出正确标签，更关心模型在不确定情形下能否保持克制、能否区分高噪声样本与低噪声样本、能否为后续人工复核或风险控制提供额外依据。若下游阶段仍然采用单一线性层或普通 MLP 对长文本特征做确定性映射，那么模型虽然可能获得较高的平均指标，却难以回答“模型为何做出该判断”以及“模型对当前样本有多大把握”这两个关键问题。

\subsection{研究问题}
基于上述背景，本文关注如下研究问题：在临床长文本已经能够通过强编码器获得高质量表示的前提下，如何设计一种兼顾判别能力、数值稳定性与不确定性表达能力的输出层，使模型不仅输出预测结果，还能够输出与样本难度相关的可信性信息，并在训练阶段利用这些信息进一步优化最终判别函数。

这一问题同时具有方法层面和工程层面的双重挑战。方法层面，临床样本之间存在明显异质性，不同病历的证据充分程度和噪声水平并不一致，因此需要一种能够显式刻画输入相关不确定性的机制。工程层面，若直接对长上下文编码器进行大规模端到端微调，不仅训练成本高，而且不利于在毕业设计周期内开展多轮实验、消融和结果分析。因此，本文更关注一种低侵入、可实现、可复用的可信建模方案，而不是重新训练一个新的医学基础模型。

\subsection{核心思路}
围绕上述问题，本文采用“冻结编码器、重构输出层”的研究路线。具体而言，首先使用 \encoder{} 作为冻结式特征提取器，将临床长文本映射为固定维度表示；其次，以线性头、MLP 头和 Ridge 头作为比较基线，建立一个从简单映射到结构化输出层的对照体系；最后提出 \method{} 框架，在特征层同时建模预测均值和输入相关方差，并在训练后期引入加权最小二乘精修机制，对最后一层参数进行结构化重估。

这一思路的关键不在于继续加深网络，而在于重新定义下游预测层的职责。传统输出头只负责把高维表示映射成标签分数，而本文提出的输出层同时承担三项功能：一是给出点预测；二是估计样本级别、标签级别的不确定性；三是利用该不确定性对训练样本进行结构化重加权。换言之，本文把长上下文临床编码器之后的下游阶段，从“普通判别层”提升为“可信输出层”。

\subsection{本文贡献}
本文的主要贡献如下。
\begin{itemize}[leftmargin=*, itemsep=2pt]
  \item 提出了一种面向临床长文本预测的可信输出层框架 \method{}。该框架以 \encoder{} 作为冻结式长上下文编码器，在特征空间中联合建模预测均值、输入相关方差和加权最小二乘精修过程，从而将深度文本表示与结构化统计求解结合起来。
  \item 给出了一套围绕可信输出层的完整研究叙事与实现路径。本文不是停留在概念层面提出“不确定性建模”的口号，而是明确给出特征提取、双分支建模、Beta-NLL 训练、WLS 精修和后续扩展接口之间的关系。
  \item 在参考性实验设置下，构建了线性头、MLP 头、Ridge 头、阶段 A 模型与阶段 B 模型之间的系统比较框架。结果趋势表明，完整 \method{} 在 weighted F1、micro F1、macro AUROC 以及误差指标上均表现出更稳定的综合优势。
  \item 将项目工程基础重组为可扩展的论文骨架，使该工作既可作为毕业设计论文，也可作为后续进一步引入 Bayesian Last Layer、保形预测和局部参数高效微调的研究起点。
\end{itemize}

\subsection{论文结构}
本文其余部分组织如下。第二节给出问题定义与任务形式化描述。第三节概述与本文相关的技术脉络。第四节详细介绍 \method{} 的模型设计、损失函数与两阶段训练策略。第五节给出实现细节与实验设置。第六节展示参考性实验结果与分析。第七节讨论方法优势、边界与适用条件。第八节总结全文并说明后续工作方向。
